%==============================================================================
% \appendix
% \section*{Open Science}
% \label{app:open_science}
% %==============================================================================

% We release LLAnalyzer together with the experimental harnesses,
% configurations, analysis scripts, patches, and supporting artifacts used
% to produce the results reported in this paper. Our goal is to make each
% boundary-specific measurement independently reproducible rather than
% provide only the final aggregate results. The artifact therefore contains
% the experimental harnesses and configuration information required to
% reproduce our model-loading, runtime-memory, wrapper-persistence, and
% serving-interface experiments. For anonymous review, the artifact is
% available at:
% \url{https://github.com/0xzodiac/Local-LLMs-Forensics}.

% Where licensing or redistribution restrictions prevent us from
% redistributing third-party model artifacts, we instead provide the exact
% model identifiers, versions, SHA-256 digests, and acquisition
% instructions. These identifiers allow an independent evaluator to
% reconstruct the model set used in our experiments and verify that the
% resulting artifacts are byte-identical to those evaluated in this study.


% %==============================================================================
% \section*{Ethical Considerations}
% \label{app:ethics}
% %==============================================================================

% We designed all experiments to avoid exposing data belonging to real
% users. We conducted the experiments exclusively on systems and accounts
% under our control and used synthetic UUID-tagged prompts as experimental
% secrets. We did not collect human-subject data, access third-party
% conversations, or test cross-tenant attacks against production systems.

% For the isolation-boundary experiments, we instantiated both the victim
% and attacker tenants ourselves. We performed state-recovery and timing
% experiments only against researcher-controlled serving instances and
% never attempted to recover prompts or conversation state from publicly
% accessible LLM servers.

% Our Internet measurements were limited to passive service metadata and
% non-destructive availability checks. We did not exploit, modify, or
% retrieve state from discovered hosts. Where our experiments uncovered
% security-relevant implementation behaviour, we followed coordinated
% vulnerability-disclosure practices and withheld artifact details where
% premature release could create unnecessary risk.


%==============================================================================
\section{Experimental Methodology}
\label{app:methodology}
%==============================================================================

We designed our experiments around a common question that applies to all
four boundaries: after a prompt crosses a particular software boundary,
can information attributable to that prompt still be recovered where the
corresponding confidentiality objective says it should no longer be
accessible?

To answer this question consistently, we use the same experimental
ground-truth strategy throughout the study. We embed fresh UUID canaries
in controlled prompts, record their association with the originating
request, execute the target operation, and then search the
boundary-specific evidence source for exact recovery. What changes
between experiments is therefore the boundary and evidence source, not
the basic criterion used to establish disclosure.

Trial counts are experiment-specific. Slot, injection, continuation, and
KV-overhead workloads use retained n=50-per-model files; the two
twelve-tenant runtime rows are single-run observations; D4 has 200 primary
records plus one separately named debug file; and Ollama is accounted for
from its retained telemetry tree. Hardware, model artifacts, inference
parameters, and software revisions are not assumed identical across these
families unless the methods below state that they are.


%------------------------------------------------------------------------------
\subsection{Hardware and Software Configuration}
\label{app:hardware}
%------------------------------------------------------------------------------

\paragraph{Primary test environment.}
We conduct the main experiments on a dedicated bare-metal Ubuntu
workstation and use the same machine for the model-file, runtime-memory,
wrapper-application, and serving-interface measurements. We avoid
unrelated user workloads during measurement runs to reduce interference
from background memory allocation, scheduling, and network activity.

Table~\ref{tab:hardware_config} records the hardware and operating-system
configuration used for the reported experiments.

\begin{table}[!b]
    \centering
    \caption{\small Hardware and operating-system configuration used in
    our experiments.}
    \label{tab:hardware_config}
    \footnotesize
    \setlength{\tabcolsep}{4pt}
    \begin{tabular}{ll}
        \toprule
        \textbf{Component} & \textbf{Configuration} \\
        \midrule
        Operating system &
        Ubuntu Linux, bare-metal 24.04.4 LTS \\
        Linux kernel &
        7.0.0-29-generic \\
        CPU &
        Intel Core i7-9850H CPU @ 2.60GHz \\
        System memory &
        16\,GB (15\,GiB reported) \\
        GPU &
        CUDA-capable GPU with 8\,GB VRAM \\
        GPU driver &
        NVIDIA 595.84 \\
        CUDA toolkit &
        12.0 (V12.0.140) \\
        File system &
        ext4 \\
        \bottomrule
    \end{tabular}
\end{table}

We run the GPU-memory experiments directly on the bare-metal Ubuntu host
rather than through a virtual machine or compatibility layer. This gives
us direct access to the GPU tooling and avoids introducing an additional
GPU-memory-management layer into the measurements. The historical
\texttt{cuda-gdb} attach outcomes and the host-side \texttt{/proc} scan are
reported separately below; direct tooling availability is not evidence that
an unprivileged process can read device memory.

\paragraph{Memory-acquisition permissions and provenance.}
The headline CPU scans use the Python collectors in
\texttt{multilayer\_v2.py}, \texttt{cross\_slot\_bleed.py}, and
\texttt{kv\_cache\_zeroize\_probe.py}. Each launches the serving process with
\texttt{subprocess.Popen}, reads writable mappings from
\texttt{/proc/<pid>/maps}, and opens/reads
\texttt{/proc/<pid>/mem}. The parent--child relationship is therefore part
of the measurement procedure under the reported Yama mode~1 setting; same
UID alone is not the claimed permission. The LM~Studio scanner records a
PID discovered by process matching rather than retaining a complete
relationship record, so its acquisition provenance is lower confidence and
is not used to broaden the headline claim.

The retained GDB investigation is a separate, authorized root-attached
diagnostic and is not evidence for the co-resident adversary model. In the
GPU campaign, \texttt{cuda-gdb --pid} attach attempts were denied, whereas
the retained host-side parent--child \texttt{/proc} scan succeeded. We therefore
claim host-memory residuals under $F$, not direct device-memory access. The
historical records do not retain complete target dumpability, capability,
user-namespace, or other-LSM metadata, so those fields remain limitations
on retrospective permission reconstruction.

\paragraph{Serving software.}
We use \texttt{llama.cpp} and its \texttt{llama-server} component for the
core runtime-memory and serving-interface experiments. For the
wrapper-level measurements, we additionally evaluate LM~Studio and
Ollama. We inspect selected components of the wider local-LLM ecosystem,
including GGUF Python tooling and SGLang, where required by the
boundary-specific experiments.

\begin{table*}[!t]
    \centering
    \caption{Software components used in our evaluation and the retained provenance status.}
    \label{tab:software_config}
    \footnotesize
    \setlength{\tabcolsep}{4pt}
    \begin{tabular}{llp{7cm}}
        \toprule
        \textbf{Component} &
        \textbf{Purpose} &
        \textbf{Version} \\
        \midrule
        \texttt{llama.cpp} &
        Runtime and serving-interface evaluation &
        Source revision 721354fbdfb7743e2be\allowbreak2183d918a3cdb9276c70f recorded; June executable not hashed or retained \\
        \texttt{llama-server} &
        Multi-slot and shared-state experiments &
        June binary linkage unknown; October CPU-only rebuild retained for reference \\
        LM~Studio &
        Wrapper logging and network behaviour &
        llmster daemon 0.0.15-2; bundled backend 2.17.0 \\
        Ollama &
        Wrapper and network behaviour &
        Version not recorded in retained run metadata \\
        SGLang &
        Reproduction check for GGUF/Jinja hardening &
        0.5.13.post1, reported historical configuration \\
        glibc allocator &
        Heap-residue and allocator experiments &
        2.39, reported historical configuration \\
        \texttt{gdb} &
        Host-memory inspection &
        15.1, reported historical configuration \\
        \texttt{cuda-gdb} &
        GPU-memory inspection &
        12.0 (V12.0.140), reported historical configuration \\
        Python &
        Measurement harnesses and analysis &
        3.14.6, reported historical configuration \\
        \bottomrule
    \end{tabular}
\end{table*}
Source revision 721354f (full SHA
721354fbdfb7743e2be\allowbreak2183d918a3cdb9276c70f) is recorded in the project
reports; the binaries executed in June 2026 were not hashed and are no
longer available. An October 2026 CPU-only rebuild is retained for
reference, is not claimed byte-identical, and its equivalence to the
executed binary is unknown. LM~Studio runs used its bundled llama.cpp
backend; the GPU runs used a CUDA build whose revision is not recorded.
Experiments were run in June 2026 on the Table~S1 hardware, are reported
from retained records, and were not re-run for this revision.

\paragraph{Runtime configurations.}
We use the default \texttt{glibc} allocator configuration as our runtime
baseline. To determine whether allocator configuration changes the
observed prompt residue, we separately evaluate:

\begin{itemize}
    \item \texttt{MALLOC\_ARENA\_MAX=1}, which restricts the number of
    per-thread allocator arenas; and
    \item \texttt{MALLOC\_PERTURB\_=165}, which overwrites allocated or
    released memory with a deterministic byte pattern.
\end{itemize}

These controls require neither source-code modification nor binary
recompilation. We evaluate them separately from the code-level
zeroisation patch so that we can distinguish allocator-level effects
from explicit application-level sanitisation. Allocator-tunable outcomes are
reported only where a numerical comparison and retained logs support them; an
unreported allocator setting is not treated as evidence of mitigation.

\paragraph{Isolation and network monitoring.}
We execute network-behaviour experiments inside an isolated network
namespace and monitor outbound connections throughout each run. We
exclude loopback communication between the wrapper and its local
inference runtime from the external-connection count.

\paragraph{Patched and unpatched binaries.}
For the runtime-mitigation experiments, we compare an unmodified baseline
against instrumented binaries implementing the sanitisation invariants
identified through our provenance analysis. The principal patch
configurations are:

\begin{description}
    \item[$I_2$:] deterministic sanitisation of released KV-cache cells;
    \item[$I_3$:] sanitisation of selected request and response buffers;
    and
    \item[$I_4$:] secure allocator handling for selected sensitive
    objects.
\end{description}

For each comparison, we keep the model, prompt sequence, inference
parameters, and runtime configuration identical between the patched and
unpatched conditions.


%------------------------------------------------------------------------------
\subsection{Experimental Corpora, Models, and Parameters}
\label{app:models_parameters}
%------------------------------------------------------------------------------

We do not use a conventional training or benchmark dataset because our
objective is not to measure model accuracy. Instead, we construct
boundary-specific experimental corpora that provide known ground truth
for each confidentiality test. These corpora comprise malformed GGUF
artifacts, synthetic UUID-tagged prompts, controlled multi-tenant
conversations, wrapper-persistence workloads, timing probes, and
Internet deployment metadata.

\paragraph{Evaluated models.}
We evaluate four model families selected to vary architecture, parameter
scale, chat-template behaviour, stopping behaviour, and reasoning style.
This diversity allows us to distinguish serving-stack behaviour that
persists across models from effects attributable to a particular model
family.

\begin{table}[!t]
    \centering
    \caption{Models used throughout the boundary-oriented
    evaluation.}
    \label{tab:evaluated_models}
    \footnotesize
    \setlength{\tabcolsep}{4pt}
    \scalebox{0.90}{\begin{tabular}{lll}
        \toprule
        \textbf{Model} &
        \textbf{Short name} &
        \textbf{Role in evaluation} \\
        \midrule
        \texttt{gemma-4-e4b-it} &
        Gemma &
        Instruction-following architecture \\
        \texttt{qwen-3.5-9b} &
        Qwen &
        Reasoning-capable model family \\
        \texttt{nemotron-3-nano-4b} &
        Nemotron &
        Compact reasoning model \\
        \texttt{phi-4-reasoning-plus} &
        Phi-4 &
        Explicit reasoning-output model \\
        \bottomrule
    \end{tabular}
}            
\end{table}

We execute all models locally using their GGUF-compatible
representations. Unless an experiment explicitly requires stochastic
generation, we use greedy decoding to reduce variation between repeated
trials.

%==============================================================================
\subsection{Confidentiality Boundaries and Adversary Models}
\label{app:boundary_summary}
%==============================================================================

Table~\ref{tab:S1} provides the comprehensive boundary-to-adversary
mapping referenced in the main manuscript.

\begin{table}[!t]
\centering
\small
\caption{Confidentiality boundaries and corresponding adversary models.}
\label{tab:S1}
\setlength{\tabcolsep}{3pt}
\begin{tabular}{p{1.7cm}p{0.7cm}p{4.5cm}}
\toprule
\textbf{Boundary} & \textbf{Adv.} & \textbf{Security objective} \\
\midrule
Integrity    & ${A}_3$ & Reject malformed model artifacts without
                         compromising trusted execution. \\
Lifetime     & $F$      & Characterize residual prompt data and conditional
                         forensic recoverability after inference. \\
Persistence  & ${A}_1$ & Prevent unintended prompt retention on local
                         storage. \\
Isolation    & ${A}_2$ & Prevent one client from recovering another
                         client's prompt or conversation state. \\
\bottomrule
\end{tabular}
\end{table}

The Isolation objective is a threat-model property. The retained slot and
continuation workloads used one unauthenticated process and did not instantiate
two authenticated clients, distinct credentials, or ownership controls.

Table~\ref{tab:S2} specifies the evidence source and
measured property for each boundary.

\begin{table}[!t]
\centering
\small
\caption{Boundary-specific evidence collected by LLAnalyzer.}
\label{tab:S2}
\begin{tabular}{p{1.6cm}p{2.2cm}p{3.1cm}}
\toprule
\textbf{Boundary} & \textbf{Evidence} & \textbf{Measured property} \\
\midrule
Integrity &
Parser execution &
Safe rejection of malformed artifacts \\

Lifetime &
Process memory &
Post-inference prompt recoverability \\

Persistence &
Wrapper artifacts &
Prompt retention on persistent storage \\

Isolation &
Serving protocol &
Shared-state separation \\
\bottomrule
\end{tabular}
\end{table}

%==============================================================================
\subsection{Model Specifications}
\label{app:model_specs}
%==============================================================================

Table~\ref{tab:S3} provides the full model specifications referenced
in the main manuscript, including parameter counts, architectures, and
artifact sizes.

\begin{table}[!th]
    \centering
    \resizebox{\columnwidth}{!}{%
        \footnotesize
        \begin{tabular}{@{}llll@{}}
            \toprule
            Model & Params & Architecture & Size (Bytes) \\
            \midrule
            NVIDIA Nemotron-3-Nano
                & 4\,B & hybrid attn/SSM & 2{,}837{,}072{,}896 \\
            Qwen3.5
                & 9\,B & transformer & 5{,}627{,}044{,}256 \\
            Gemma-4-E4B-it
                & 7.5\,B & transformer & 5{,}335{,}289{,}664 \\
            Phi-4-reasoning-plus
                & 14\,B & transformer (reasoning) & 9{,}053{,}116{,}128 \\
            \bottomrule
        \end{tabular}%
    }
    \caption{Open-weight models evaluated by LLAnalyzer. All model 
    artifacts use GGUF Q4\_K\_M quantisation and were obtained from the 
    \texttt{lmstudio-community} organization on Hugging Face.}
    \label{tab:S3}
\end{table}

\paragraph{Model-file corpus.}
We evaluate the integrity boundary using complementary structured and
coverage-guided approaches. The structured corpus contains malformed GGUF
artifacts targeting tensor dimensions, names, metadata lengths, header
fields, and related consistency constraints. We additionally exercise
the parser through coverage-guided fuzzing to explore states not captured
by manually constructed mutations.

For the structured component, we use:

\begin{enumerate}
    \item an 80-case malformed-input battery targeting the
    post-allocation-validation path of the C++ GGUF reader; and
    \item a 21-file structural-header corpus containing malformed tensor
    dimensions, tensor names, metadata lengths, and related header
    inconsistencies.
\end{enumerate}

We load each artifact independently and classify the resulting behaviour
as clean rejection, handled parsing error, unhandled exception, process
termination, successful malformed loading, or memory-safety failure. We
also process the 21 structural-header artifacts through the repository's
\texttt{gguf-py} tooling to compare behaviour between the C++ and Python
parsers.

Older manuscript text described a 24-hour AFL++ campaign, but no campaign
log, queue, executable, coverage record, or manifest is retained. This
revision does not use that campaign total as evidence.

\paragraph{Synthetic confidentiality workload.}
We establish experimental ground truth using fresh UUID canaries
generated separately for each experimental client and trial. Before injecting a
canary, we verify that it does not occur in the corresponding
pre-experiment evidence source. We then embed the identifier directly
into the controlled prompt and record its association with the
originating request.

This design gives us a direct attribution criterion. When $F$
subsequently recovers an exact canary from process memory, a
wrapper-managed file, a protocol response, or a continuation response, we
can associate that occurrence with the request in which we introduced it
rather than infer its provenance from surrounding plaintext. The phrase
``another client's execution context'' is a threat-model possibility, not a
capability exercised by the retained slot or continuation workloads: those
workloads used one unauthenticated process and did not provide distinct
clients, credentials, or ownership controls.

Depending on the confidentiality boundary under evaluation, we record a
direct disclosure when the originating request's canary appears in:

\begin{itemize}
    \item process heap memory collected through the authorized $F$ path;
    \item KV-cache or associated runtime buffers;
    \item GPU or pinned host memory;
    \item wrapper-generated files or logs;
    \item protocol response content;
    \item reasoning-output fields; or
    \item another client's continuation response.
\end{itemize}

\paragraph{Runtime-memory workload.}
The retained headline twelve-tenant scans are two Nemotron single-run
observations. The separate multilayer-v2 diagnostic family is Phi-4-only,
while the allocator probe contains paths for all four models but retains
only aggregate summary evidence rather than per-repetition records. The
four-model slot, continuation, injection, KV-overhead, and D4 families are
described separately below and should not be conflated with the
twelve-tenant rows. For each retained workload, $F$ is the harness process
that launches the serving process as a direct child. The Python scanners
open \texttt{/proc/<pid>/maps} and then read writable regions through
\texttt{/proc/<pid>/mem}; these reads are subject to Linux's
\texttt{PTRACE\_MODE\_ATTACH\_FSCREDS} checks~\cite{linux_process_permissions}.
The retained runtime-memory protocols include:

\begin{enumerate}
    \item the unmodified runtime;
    \item each code-level mitigation configuration;
    \item \texttt{MALLOC\_ARENA\_MAX=1};
    \item \texttt{MALLOC\_PERTURB\_=165}; and
    \item selected combinations of code-level and allocator-level
    mitigations.
\end{enumerate}

We measure both the number of residual cleartext prompt copies and the
aggregate number of bytes associated with matching regions. The KV-cache
diagnostic code can enumerate 56 targeted cache cells after release, but no
per-cell result file is retained. We therefore treat that number as a
configured diagnostic scope, not as an independently verified headline
outcome.

\paragraph{Wrapper-application workload.}
We submit repeated UUID-tagged prompts to LM~Studio and Ollama and inspect
their application-managed files before and after each run. We evaluate
the reported logging configuration and, where supported, repeat the
measurement after explicitly disabling sensitive-data logging. For
LM~Studio, the retained mitigation report identifies llmster daemon
0.0.15-2 and bundled backend 2.17.0. The original plaintext log is not
retained, and no clean-install default is claimed.

The retained Ollama tree contains 206 metadata records, 203 successful
trigger files, two explicit model-memory failures, and one metadata-only
incomplete record. No Ollama executable version is recorded in those
metadata files. Six LM~Studio hook summaries show external connections but
no canary in visible captured or hooked payloads; TLS was not decrypted.

\paragraph{Shared-process slot workload.}
The retained slot script uses one unauthenticated loopback server and one
Python process. It creates a UUID-bearing request in slot~0, saves an
attacker-selected filename, restores that file into slot~1, and scores
success only when \texttt{n\_saved = n\_restored}. No second client,
distinct API key, wrong-key control, or ownership check was tested.

We retain 200 slot-transfer equality trials in total, comprising 50 per
model. Each trial consists of:

\begin{enumerate}
    \item creating a victim session containing a unique secret;
    \item saving the victim's slot state;
    \item restoring that file into slot~1; and
    \item comparing the returned saved and restored token counts.
\end{enumerate}

We separately evaluate model-level continuation. The retained files are
dated July 3--4, 2026 and contain 126/200 KV-channel hits (Gemma 50/50,
Qwen 50/50, Nemotron 25/50, Phi-4 1/50) and 200/200 memory-channel hits.
These records are not attributed here to the June campaign and do not test
distinct credentials. The logical protocol is shared, but retained GPU
offload settings differ by model: 999 layers for Gemma, Qwen, and Nemotron,
and 20 layers for Phi-4. This configuration difference is retained as a
nuisance factor rather than treated as a balanced hardware comparison.

\paragraph{Cross-model residue workload (D4).}
The retained D4 family contains 200 timestamped primary records, divided
into 50 Gemma-to-Phi-4, 50 Qwen-to-Gemma, 50 Nemotron-to-Qwen, and 50
Phi-4-to-Nemotron swaps. Each trial keeps the same serving process while
model~A is replaced by model~B, then scans writable mappings for the
model-A canary. All 200 records retain the same process identity and a
positive post-swap model-A memory observation. One separately named debug
JSON is excluded from the primary denominator. The measurement is a
byte-occurrence residue observation, not a model-output disclosure or a
two-principal authorization test.

\paragraph{Timing workload.}
An older report describes a cached-versus-uncached TTFT workload, but the
per-request vectors and original analysis script are not retained. We
therefore treat timing separation as historical context only and do not use
it as an independently validated quantitative result or deployment-wide
claim.

\paragraph{Reasoning-channel workload.}
We separately inspect conventional response content and
reasoning-specific response fields when evaluating model-dependent
disclosure. We record a disclosure when the target secret appears in
either channel. We make this distinction because some evaluated models,
particularly Qwen and Phi-4, may place recovered information in
reasoning fields that are not exposed through the conventional response
content.

\paragraph{External-scope limitation.}
No Internet-prevalence count is used as evidence in this revision. All
shared-serving observations are from researcher-controlled processes and
synthetic canaries. No vendor decision or finding-specific disposition is
inferred from policy text.

\paragraph{Disclosure status.}
No disclosure report has been transmitted and no vendor decision is
established for the observations reported here. The dated
\texttt{DISCLOSURE\_PLAN.md} entry is a researcher-side scope determination.
Policy text is not treated as a finding-specific assessment. We present the
saved-state and wrapper observations as controlled findings and
deployment-hardening guidance.

\paragraph{Summary of experimental scales.}
Table~\ref{tab:experiment_scale} summarises the principal workloads used
throughout the study.

\begin{table}[!t]
    \centering
    \caption{Scale of the principal experimental workloads.}
    \label{tab:experiment_scale}
    \footnotesize
    \setlength{\tabcolsep}{4pt}
    \resizebox{\columnwidth}{!}{\begin{tabular}{lrl}
        \toprule
        \textbf{Experiment} &
        \textbf{Scale} &
        \textbf{Unit} \\
        \midrule
        Structured GGUF battery &
        80 &
        malformed inputs \\
        Structural-header analysis &
        21 &
        malformed files \\
        Coverage-guided fuzzing &
        Not retained &
        historical campaign total \\
        KV-cell sanitisation &
        56 &
        scope only; per-cell records absent \\
        Same-process slot transfer &
        200 &
        trials; 50 per model \\
        KV continuation &
        200 &
        trials; 50 per model \\
        Timing side channel &
        Historical only &
        vectors/script not retained \\
        Ollama telemetry &
        206 metadata; 203/2/1 &
        successful/failed/incomplete \\
        \bottomrule
    \end{tabular}
}
\end{table}


%==============================================================================
\section{Statistical Analysis}
\label{app:statistics}
%==============================================================================

We choose the statistical treatment according to the observable produced
by each boundary experiment rather than applying a single statistical
test across all measurements. Our analysis emphasises repeated
measurement, effect magnitude, and explicit reporting of both positive
and negative results.

\paragraph{Binary outcomes.}
For experiments producing a success-or-failure outcome, we report the
observed proportion

\[
\hat{p}=\frac{x}{n},
\]

where \(x\) is the number of successful disclosures or anomalous trials
and \(n\) is the total number of trials.

We report binomial uncertainty using two-sided 95\% Wilson score
intervals~\cite{wilson1927}. We use Wilson intervals because several of
our measurements lie close to the boundaries of zero or one, where
normal approximations are inappropriate. For example, observing zero
successes in 50 trials does not establish a true probability of zero;
the corresponding Wilson interval instead provides an explicit bound on
the unobserved event probability.

\paragraph{Deterministic findings.}
We use the term \emph{deterministic} only when every repeated trial under
a fixed experimental condition produces the same outcome. Examples
include direct slot-state equality succeeding in 200 of 200 trials. Our use of
the term refers to repeatability under the evaluated configuration and
does not imply universal behaviour across all software versions,
hardware platforms, or deployment environments.

\paragraph{Memory-residue measurements.}
For runtime-memory experiments, we report:

\begin{itemize}
    \item the number of residual cleartext copies;
    \item the aggregate bytes occupied by matching memory regions; and
    \item the direct observed values for each retained configuration.
\end{itemize}

The two Nemotron twelve-tenant configurations are single-run observations with
different executable and runtime settings. We therefore report their copy
and byte metrics directly and do not calculate a causal relative-reduction
effect. Copy counts are byte occurrences in scanned mappings, not
independent secrets, tenants, or trials.

\paragraph{Timing-side-channel analysis.}
The historical timing vectors and analysis script are not retained. We do
not recompute or validate the historical timing separation in this revision,
and no timing number is used in the main quantitative conclusions.

\paragraph{Cross-model reporting.}
We report results separately for Gemma, Qwen, Nemotron, and Phi-4 where
model-specific behaviour affects the outcome. We avoid pooling these
measurements into a single aggregate proportion because differences in
chat templates, stopping behaviour, reasoning channels, and allocation
patterns are themselves relevant to the confidentiality analysis.

We describe a result as applying across model families only for a named
retained four-model family (slot equality, July continuation, direct
injection, KV overhead, D4 residue, LM~Studio hooks, or Ollama telemetry).
The Nemotron-only twelve-tenant and GPU-v5 rows, Phi-4-only multilayer
diagnostics, and aggregate-only allocator summary are not generalized to
the other models. Where quantitative differences occur, we retain and
discuss them rather than obscuring them through pooled averages.

\paragraph{Analysis of null results.}
We retain negative results when they meaningfully constrain a plausible
security hypothesis rather than omitting experiments that fail to
produce an attack. Before interpreting an observation as a negative
result, we verify that:

\begin{enumerate}
    \item the relevant instrumentation and response fields are active;
    \item an appropriate positive-control condition produces the expected
    observable;
    \item the serving process remains operational after the trial; and
    \item the null outcome is reproducible under the stated experimental
    condition.
\end{enumerate}

We do not infer universal security from any null result. Negative
observations are retained only with the instrumentation, controls, and
deployment conditions that produced them.

\paragraph{Reproducibility.}
For each quantitative claim, we identify the retained measurement script,
configuration, raw experimental log, or historical report in the audit
artifacts. The harness generates randomised identifiers and UUID canaries
and records them alongside the trial in which they are introduced. Claims
whose historical per-request data or binary linkage are not retained are
marked as such rather than described as independently recomputed.
% 	\section*{Appendix}
% 	\renewcommand{\thesubsection}{\Alph{subsection}}
% 	\setcounter{subsection}{0}
	
% 	\subsection{Probabilistic Reconstruction of the KV-Continuation
% 		Channel --- A Null Result that Bounds the Per-Model Rates
% 		\label{app:B}}
% 	The KV-continuation channel of \S\ref{sec:scope} recovers the
% 	sender's secret via the greedy (argmax) decode path after a
% 	cross-tenant \texttt{/slots} restore. The per-model rates are
% 	gemma-4-e4b-it 50/50, qwen-3.5-9b 50/50, nemotron-3-nano-4b
% 	21/50, and phi-4-reasoning-plus 3/50. A natural objection to
% 	the phi-4 3\% rate is that the greedy path halted at
% 	end-of-sequence before the secret's tokens, while probability
% 	mass for those tokens remained in the top-$K$ distribution.
% 	Primitive~B (\emph{probabilistic reconstruction}) sets
% 	$n_{\mathrm{probs}}\!=\!20$ and scores a hit if the secret's
% 	leading eight characters appear anywhere in the top-20 token
% 	string across all harvested positions, not only as the greedy
% 	choice. Per-model results at $n{=}50$:
	
% 	\begin{table}[!t]
% 		\centering
% 		\begin{tabular}{lrrr}
% 			\toprule
% 			Model & Greedy hits & Prob hits & Prob recovers \\
% 			\midrule
% 			\texttt{gemma-4-e4b-it}        & 50/50 & 50/50 & 0 \\
% 			\texttt{qwen-3.5-9b}          & 50/50 & 50/50 & 0 \\
% 			\texttt{nemotron-3-nano-4b}   & 21/50 & 21/50 & 0 \\
% 			\texttt{phi-4-reasoning-plus}  &  3/50 &  3/50 & 0 \\
% 			\bottomrule
% 		\end{tabular}
% 		\caption{Probabilistic reconstruction of the KV-continuation
% 			channel. The top-20 distribution never recovers a greedy
% 			miss. The reported per-model rates in
% 			\S\ref{sec:scope} are a true ceiling on the channel
% 			under both argmax and top-$K$ decoding, not a
% 			decoding-strategy artifact.}
% 		\label{tab:primitiveB}
% 	\end{table}
	
% 	\noindent\textbf{Why the null result is real, not a failure of the
% 		experiment.} We verified (i) the top-20 distribution is being
% 	harvested correctly by inspecting the raw response schema
% 	(\texttt{completion\_probabilities} with
% 	\texttt{top\_probs[]} per position); (ii) on the models where
% 	the greedy channel already succeeds (gemma-4, qwen-3.5), the
% 	probabilistic channel matches the greedy channel exactly
% 	($p_{\mathrm{prob}}\!=\!p_{\mathrm{greedy}}$), confirming the
% 	parse is identical between success and miss cases; and
% 	(iii) the miss cases on phi-4 and the misses on nemotron are
% 	not in the top-20 --- the secret's tokens carry no
% 	non-trivial probability mass at the harvested positions.
% 	This confirms the per-model rates in \S\ref{sec:scope} are a
% 	genuine upper bound on the channel under any single-prompt
% 	decode, not a greedy-decoding artifact.
	
% 	\subsection{Stealth KV Backdoor --- A Null Result that Bounds the
% 		Integrity Threat to Confidentiality-Only
% 		\label{app:D}}
% 	The cross-tenant \texttt{/slots} primitive transfers KV state
% 	between slots; Primitive~D asks whether an attacker who
% 	restores a \emph{crafted} KV state into a victim's slot can
% 	make the victim's own subsequent generation follow a
% 	hidden instruction the victim's visible prompt never
% 	contains --- an integrity attack, not a confidentiality leak.
% 	The benign proof-of-control is a marker
% 	\texttt{[SILENTSLOT-OWNED-<id>]} that the injected state
% 	instructs the victim model to prepend. We restore the
% 	crafted state, then issue a benign unrelated question to
% 	the same slot, and score a hit if the marker appears in
% 	the response. Per-model results at $n{=}50$:
	
% 	\begin{table}[!t]
% 		\centering
% 		\begin{tabular}{lr}
% 			\toprule
% 			Model & Behavioral control achieved \\
% 			\midrule
% 			\texttt{gemma-4-e4b-it}        & 0/50 \\
% 			\texttt{qwen-3.5-9b}          & 0/50 \\
% 			\texttt{nemotron-3-nano-4b}   & 0/50 \\
% 			\texttt{phi-4-reasoning-plus}  & 0/50 \\
% 			\bottomrule
% 		\end{tabular}
% 		\caption{Stealth KV backdoor at $n{=}50$. Zero control
% 			achieved on any model. Wilson 95\% CI on every model
% 			upper-bounds at $0.07$. The cross-tenant
% 			\texttt{/slots} primitive is a confidentiality attack
% 			on the same-channel KV state, not a cross-channel
% 			behavioral-injection vector.}
% 		\label{tab:primitiveD}
% 	\end{table}
	
% 	\noindent\textbf{Why the null result is real.} The same
% 	prefix-reset mechanism that caused the original 0/13
% 	KV-continuation result (and was fixed for the \emph{read}
% 	direction via the \texttt{/apply-template} bypass in
% 	\S\ref{sec:scope}) defeats the \emph{write} direction here:
% 	the victim's chat-template rendering, applied to its own
% 	benign question, discards the non-shared backdoor prefix
% 	before generation, so the injected instruction never
% 	conditions the output. This is the same null class as the
% 	one the original P1 pilot demonstrated for the read
% 	direction; the fix does not transfer because the
% 	write-path bypass would require the victim to issue the
% 	\texttt{/apply-template} request, which it does not.
% 	The honest framing is that the threat is bounded to
% 	\emph{confidentiality on the same-channel KV state}; a
% 	successful integrity injection would require a different
% 	vector (e.g. poison the chat template at the \texttt{tokenizer.chat\_template}
% 	field, which the paper's existing
% 	\S\ref{sec:rel-gguf} ecosystem analysis already covers at
% 	the file-format layer).
	
% 	\subsection{Concurrency Race (TOCTOU) on \texttt{/slots} Save/Restore ---
% 		A Null Result that Bounds the Race-Surface
% 		\label{app:E}}
% 	A TOCTOU race is a bug
% 	where the state a request validated changes before the request acts on it. This appendix tests for one.
% 	Primitive~E fires four concurrent slot operations per trial
% 	(save-slot-0, restore-slot-1-from-slot-0, save-slot-1,
% 	restore-slot-0-from-slot-1) to maximize interleaving of the
% 	checkpoint-search path. The post-race state is checked
% 	for HTTP errors, server death, and cross-tenant state
% 	contamination (a slot echoing the other slot's seed
% 	secret). Per-model results at $n{=}50$:
	
% 	\begin{table}[!t]
% \centering
% \caption{Concurrency TOCTOU on \texttt{/slots} save/restore at
% $n{=}50$. Anomaly trials (AT) report HTTP-level request
% failures (recoverable errors; the server remains
% healthy and serves subsequent requests). Server crashes
% (SC) and cross-tenant contamination (CTC) are zero
% across all four models. The save/restore path is
% isolation-safe under the tested race workload.}
% \label{tab:primitiveE}
% \setlength{\tabcolsep}{4pt}
% \begin{tabular*}{\columnwidth}{@{\extracolsep{\fill}}lccc}
% \toprule
% Model & AT & SC & CTC \\
% \midrule
% \texttt{gemma-4-e4b-it}       & 20/50 & 0 & 0 \\
% \texttt{qwen-3.5-9b}          & 20/50 & 0 & 0 \\
% \texttt{nemotron-3-nano-4b}   & 17/50 & 0 & 0 \\
% \texttt{phi-4-reasoning-plus} & 26/50 & 0 & 0 \\
% \bottomrule
% \end{tabular*}
% \end{table}
	
% 	\noindent\textbf{Why the null result is real.} The HTTP errors
% 	observed are recoverable (the server continues to serve
% 	on the same socket), and the cross-tenant check after
% 	the race is unambiguous: each slot's post-race generation
% 	echoes only the seed secret of the request that completed
% 	on that slot. The race-safe negative is itself a
% 	publication-quality result --- it bounds the
% 	multi-tenant race surface for this primitive and shows
% 	that an attacker who relies on race timing gains nothing
% 	beyond occasional request errors (visible to both
% 	parties, recoverable by the client).
	
% 	\subsection{Total-Surface-Decomposition Audit
% 		\label{app:audit}}
% 	Beyond the four primitives above, we conducted a
% 	\emph{parallel-code-audit pass} of the \texttt{llama.cpp}
% 	server tree, with each of four sonnet subagents owning one
% 	unexplored surface. Every agent was scoped to read-only
% 	analysis of the local source tree and to design a
% 	localhost-only experiment for the surface it owned; no
% 	external hosts were scanned or probed. The findings:
	
% 	\begin{description}
% 		\item[\textbf{HTTP-endpoint surface, $\bm{0}$ (clean
% 			negative).}] \texttt{/infill}, \texttt{/embedding},
% 		\texttt{/rerank}, \texttt{/tokenize}, and \texttt{/props}
% 		expose no per-tenant state and are not cross-tenant
% 		vectors. \texttt{/rerank} uses a
% 		\texttt{string\_replace\_all} substitution on the model's
% 		\texttt{rerank} chat template, \emph{not} a Jinja render;
% 		it is not in an unrelated template-escape class. The
% 		process-global \texttt{/lora-adapters} channel is
% 		documented as a separate preliminary finding in
% 		Appendix~\ref{app:LoRA}.
% 		\item[\textbf{Minja chat-template SSTI, $\bm{0}$ (clean
% 			negative).}] This build ships its own in-tree jinja
% 		engine (\texttt{common/jinja/}, a Jinja2 subset) that
% 		replaces the upstream \texttt{minja} library. The engine
% 		sandbox-checks attribute access, and the attack class of
% 		SGLang's unrelated template-escape pattern
% 		does not parse: every attribute access through
% 		\texttt{.\,} on a non-builtin object is rejected. A
% 		malicious GGUF's \texttt{tokenizer.chat\_template} cannot
% 		achieve SSTI when the server loads it with
% 		\texttt{--jinja}.
% 		\item[\textbf{Slot-restore blob parser, $\bm{0}$ (clean
% 			negative).}] The deserializer at
% 		\texttt{llama-context.cpp} / \texttt{llama-kv-cache.cpp}
% 		validates every attacker-controlled length field read
% 		from the file (\texttt{cell\_count}, \texttt{n\_seq\_id},
% 		\texttt{v\_trans}, \texttt{n\_layer}, key types and
% 		row sizes) against the live model's
% 		\texttt{layers.size()} and the per-row size computed
% 		from \texttt{ggml\_row\_size} before any allocation or
% 		memcpy. The single-seq path rejects
% 		\texttt{n\_seq\_id $\neq$ 1}. We did not find a
% 		memory-safety candidate in the audited parse
% 		sequence.
% 		\item[\textbf{Sampler/grammar cross-slot bleed, $\bm{0}$
% 			(clean negative).}] \texttt{slot.smpl.reset(\\
% 			common\_sampler\_init(model\_tgt, task.params.sampling))}
% 		inside \texttt{launch\_slot\_task}
% 		(server-context.cpp:~1824) and
% 		\texttt{common\_sampler\_reset} (line:~369) on
% 		\texttt{init\_sampler} clear all stateful-sampler
% 		residual state, and \texttt{reset()} wipes
% 		\texttt{json\_schema} (line:~339) and calls
% 		\texttt{llama\_set\_sampler(ctx\_tgt, id, nullptr)}
% 		(line:~359). Under the canonical
% 		tenant-A-completes~$\to$~tenant-B-starts path, no
% 		observable bleed. The design is defensively correct.
% 	\end{description}
	
% 	\noindent The audit's \emph{net} contribution to the paper is a
% 	bound on four adjacent surfaces and one well-scoped
% 	follow-up experiment (the LoRA behavioural-shift
% 	reproduction, to be reported separately). The body
% 	findings (\S\ref{sec:r-protocol}--\ref{sec:r-bypass})
% 	together with this appendix establish that the
% 	cross-tenant threat in the consumer local-LLM stack is
% 	\emph{real, broad, and partially orthogonal to the
% 		flag-level mitigations} the vendor already documents, and
% 	\emph{bounded} on the surfaces where intuition would
% 	predict a worse outcome.
	
% 	\subsection{Preliminary Finding --- \texttt{/lora-adapters} Process-Global
% 		Mutation Crosses Tenants
% 		\label{app:LoRA}}
% 	The HTTP-endpoint audit of Appendix~\ref{app:audit} identified
% 	a \emph{control-channel} cross-tenant primitive: the
% 	\texttt{POST /lora-adapters} handler mutates
% 	\texttt{params\_base.lora\_adapters} (server-context.cpp,~line~2723),
% 	which is a \emph{process-global} handle. The maintainer's own
% 	\texttt{TODO @ngxson : make lora\_adapters a dedicated member of
% 		server\_context} at the same lines is the upstream self-flag.
% 	A subsequent tenant request that does not carry a
% 	per-request \texttt{lora} field silently inherits the swapped
% 	adapter via \texttt{slot.lora = params\_base.lora\_adapters}
% 	at server-context.cpp,~line~1763. On a multi-tenant server
% 	sharing one \texttt{--api-key}, the \texttt{POST} handler is
% 	gated by the same authentication as the rest of the API but
% 	not by any per-tenant scope check, and the mutation persists
% 	for the lifetime of the process.
	
% 	This is recorded as a \emph{preliminary} finding because the
% 	end-to-end behavioural-shift experiment requires (i)~a
% 	clean \texttt{llama-server} binary independent of the
% 	paper's \texttt{secure\_zero(res->content)} patch, and (ii)~a
% 	real LoRA adapter for one of the four test models. The
% 	harness we built (\texttt{scripts/silentslot\_LoRA\_hot\_swap.py})
% 	ran on the only binary available in the local tree, which is
% 	the patched one; the behavioural-shift measurement was
% 	confounded by the patch's own response-zeroisation bug
% 	(Appendix~\ref{app:audit},~\S\ref{sec:r-bypass}). The
% 	\emph{control-channel} result --- that the global is reachable,
% 	mutation persists, and the slot's fallback to the global is
% 	unconditional on the absence of a per-request \texttt{lora}
% 	field --- is verified by the source-of-truth citations
% 	above and is itself a publication-quality audit finding.
	
% 	Reproducing the full behavioural-shift result on a clean
% 	binary is a clear next step; we plan to include it in a
% 	revised version of the paper once the cross-binary
% 	experiment is complete.
	
% 	\subsection{Structural-Header Fuzzing: C++ Reader Hardened,
% 		\texttt{gguf-py} Tooling Is Not
% 		\label{app:gguf-header}}
% 	\S\ref{sec:r-file} reports 21 malformed GGUF headers rejected
% 	cleanly by the C++ reader. The same 21-file corpus was also loaded
% 	through \texttt{gguf-py}'s \texttt{gguf\_dump} utility --- the
% 	Python reference parser shipped in the same \texttt{llama.cpp}
% 	repository and used by third-party conversion, quantisation, and
% 	metadata-inspection tooling that does not link against
% 	\texttt{ggml/src/gguf.cpp} directly. Unlike the C++ reader, several
% 	cases (\texttt{zero\_tensor\_dimension}, \texttt{oversized\_tensor\_name},
% 	\texttt{gigantic\_tensor\_dim}) do not produce a caught,
% 	user-facing error: they raise an \emph{unhandled} Python exception
% 	(\texttt{IndexError} / \texttt{struct.error}) from inside
% 	\texttt{gguf\_reader.py}'s \texttt{\_get\_field\_parts}, terminating
% 	the calling process with a raw traceback rather than a diagnostic
% 	message.
	
% 	\begin{findingbox}
% 		0/21 process crashes in either parser (no memory corruption). But
% 		the C++ reader's bounds-checked design (\S\ref{sec:r-file}) is not
% 		mirrored in \texttt{gguf-py}: a subset of the same malformed
% 		headers that the C++ reader rejects gracefully instead crash the
% 		Python tooling with an unhandled exception. Any pipeline that calls
% 		\texttt{gguf-py} programmatically on an untrusted file (automated
% 		conversion services, metadata scrapers, HuggingFace-hub ingestion
% 		scripts) without wrapping the call in its own exception handling
% 		inherits an uncaught-exception denial-of-service on attacker-supplied
% 		input, despite the underlying C++ server being hardened against the
% 		identical byte patterns.
% 	\end{findingbox}
	
% 	This is a tooling-hardening gap, not a memory-safety vulnerability:
% 	impact is bounded to process termination of the calling script, with
% 	no evidence of memory corruption or code execution. We flag it
% 	because it is a natural, low-effort target for an operator who
% 	assumes ``the parser is hardened'' from the C++ side and pipes
% 	untrusted files through the Python tooling on that assumption alone.
	
% 	%==============================================================================
% 	% USENIX Open Artifacts / Artifacts Available statement
% 	%==============================================================================
% 	\section*{Open Science and Artifacts}
	
% 	The complete measurement harness, per-phase logs, the
% 	$n\!\ge\!50$ / 4-of-4-model statistical-power audit results, runtime
% 	zeroisation patches, the cross-tenant exploit re-implementation, the
% 	memorisation-volume extraction script, the network-namespaces
% 	isolation harness, and the responsible-disclosure record are
% 	released publicly under the MIT License (code) and CC-BY (data) at:
% 	\begin{center}
% 		\url{https://github.com/youseefhamdi/Local-LLM-Forensics}
% 	\end{center}
% 	The package is additionally archived via Zenodo to ensure long-term
% 	accessibility and citability:
% 	\begin{center}
% 		\url{https://doi.org/10.5281/zenodo.20998459}
% 	\end{center}
	
% 	%==============================================================================
% 	% CRediT authorship contribution statement
% 	%==============================================================================
% 	\section*{Authorship Contributions}
	
% 	\textbf{Youssef Hamdi Zafaan Ibrahim:}
% 	Conceptualisation, Methodology, Software, Investigation, Data Curation,
% 	Formal Analysis, Resources, Validation, Visualisation, Writing ---
% 	Original Draft, Writing --- Review \& Editing.
	
% 	\textbf{Mohammed Khalaf Salama:}
% 	Conceptualisation, Methodology, Investigation, Formal Analysis,
% 	Writing --- Review \& Editing.
	
% 	%==============================================================================
% 	% Declaration of Generative AI use in writing process
% 	%==============================================================================
% 	\section*{Declaration of Generative AI and AI-Assisted Technologies}
	
% 	During the preparation of this manuscript, the authors used
% 	AI-assisted language tools for grammar correction and formatting
% 	assistance. All scientific content---including the four-layer
% 	measurement design, file-format fuzzing battery, heap-residue
% 	forensics methodology, live-debugger root-cause analysis,
% 	cross-tenant exploit re-implementation, memorisation volume
% 	extraction, network-namespaces isolation harness, the patch stack
% 	with formal invariants $I_2$/$I_3$/$I_4$, and all numerical
% 	results---was conceived, executed, and verified by the authors.
% 	Every quantitative claim in this paper is supported by a released
% 	artefact: the measurement harness, per-phase logs, the
% 	$n\!\ge\!50$ / 4-of-4-model statistical-power audit results, runtime
% 	patches, and the responsible-disclosure record. The authors reviewed
% 	and edited all AI-assisted content and take full responsibility for
% 	the integrity of the published work.
	
% 	%==============================================================================
% 	% Disclosure of any competing interests
% 	%==============================================================================
% 	\section*{Disclosure of Competing Interests}
	
% 	The authors declare that they have no known competing financial
% 	interests or personal relationships that could have appeared to
% 	influence the work reported in this paper.
	
